\begin{thebibliography}{99}
\bibitem{AHJ} O. Angel, T. Hutchcroft \& A. Járai - “On the tail of the branching random walk local time” , Probab. Theory Related Fields 180 (2021) no. 1-2, p. 467-494
\bibitem{AN} K. B. Athreya \& P. E. Ney - Branching processes , Dover Publications, Inc., Mineola, NY, 2004, Reprint of the 1972 original
\bibitem{ASS23} A. Asselah, B. Schapira \& P. Sousi - “Local times and capacity for transient branching random walks” , 2023
\bibitem{BC12} I. Benjamini \& N. Curien - “Recurrence of the $\mathbb{Z}^d$-valued infinite snake via unimodularity” , Electron. Comm. Probab. 17 (2012), article ID 1, 10 pages
\bibitem{BH21} T. Bai \& Y. Hu - “Capacity of the range of branching random walks in low dimensions” , Trudy Mat. Inst. Steklov. 316 (2022), p. 32-46
\bibitem{BH22} T. Bai \& Y. Hu - “Convergence in law for the capacity of the range of a critical branching random walk” , Ann. Appl. Probab. 33 (2023) no. 6A, p. 4964-4994
\bibitem{BHJ23} N. Berestycki, T. Hutchcroft \& A. Jego - “Thick points of 4D critical branching Brownian motion” , 2023
\bibitem{BW20} T. Bai \& Y. Wan - “Capacity of the range of tree-indexed random walk” , Ann. Appl. Probab. 32 (2022) no. 3, p. 1557-1589
\bibitem{DE} A. Dvoretzky \& P. Erdős - “Some problems on random walk in space” , in Proc. Second Berkeley Symposium on Mathematical Statistics and Probability, 1950 , Univ. California Press, Berkeley-Los Angeles, CA, 1951, p. 353-367
\bibitem{DKLT} T. Duquesne, R. Khanfir, S. Lin \& N. Torri - “Scaling limits of tree-valued branching random walks” , Electron. J. Probab. 27 (2022), article ID 16, 54 pages
\bibitem{Kes} H. Kesten - “Branching random walk with a critical branching part” , J. Theoret. Probab. 8 (1995) no. 4, p. 921-962
\bibitem{LGL1} J.-F. Le Gall \& S. Lin - “The range of tree-indexed random walk in low dimensions” , Ann. Probab. 43 (2015) no. 5, p. 2701-2728
\bibitem{LGL2} J.-F. Le Gall \& S. Lin - “The range of tree-indexed random walk” , J. Inst. Math. Jussieu 15 (2016) no. 2, p. 271-317
\bibitem{LL10} G. F. Lawler \& V. Limic - Random walk: a modern introduction , Cambridge Studies in Advanced Math. , vol. 123 , Cambridge University Press, Cambridge, 2010
\bibitem{LSS24} A. Legrand, C. Sabot \& B. Schapira - “Recurrence and transience of the critical random walk snake in random conductances” , 2024
\bibitem{LZ11} S. P. Lalley \& X. Zheng - “Occupation statistics of critical branching random walks in two or higher dimensions” , Ann. Probab. 39 (2011) no. 1, p. 327-368
\bibitem{NV75} S. V. Nagaev \& N. V. Vahrušev - “An estimate of large deviation probabilities for a critical Galton-Watson process” , Teor. Verojatnost. i Primenen. 20 (1975), p. 181-182
\bibitem{NV03} S. V. Nagaev \& V. I. Vakhtel - “Limit theorems for probabilities of large deviations of a Galton-Watson process” , Diskret. Mat. 15 (2003) no. 1, p. 3-27, translation in Discrete Math. Appl. 13 (2003), no. 1, p. 1–26
\bibitem{PZ19} E. B. Procaccia \& Y. Zhang - “Connectivity properties of branching interlacements” , ALEA Lat. Am. J. Probab. Math. Stat. 16 (2019) no. 1, p. 279-314
\bibitem{Shi} Z. Shi - Branching random walks , Lect. Notes in Math. , vol. 2151 , Springer, Cham, 2015
\bibitem{Zhu1} Q. Zhu - “On the critical branching random walk I: branching capacity and visiting probability” , 2016
\bibitem{Zhu2} Q. Zhu - “On the critical branching random walk II: branching capacity and branching recurrence” , 2016
\bibitem{Zhu18} Q. Zhu - “Branching interlacements and tree-indexed random walks in tori” , 2018
\bibitem{Zhu4} Q. Zhu - “An upper bound for the probability of visiting a distant point by a critical branching random walk in $\mathbb{Z}^4$” , Electron. Comm. Probab. 24 (2019), article ID 32, 6 pages
\bibitem{Zhu3} Q. Zhu - “On the critical branching random walk III: The critical dimension” , Ann. Inst. H. Poincaré Probab. Statist. 57 (2021) no. 1, p. 73-93
\end{thebibliography}
